\documentclass[11pt,a4paper]{article}

\usepackage[margin=1in]{geometry}
\usepackage{amsmath,amssymb,amsthm}
\usepackage{graphicx}
\usepackage{booktabs}
\usepackage{microtype}
\usepackage{xcolor}
\usepackage[hidelinks,colorlinks=true,linkcolor=blue!50!black,citecolor=blue!40!black,urlcolor=blue!50!black]{hyperref}
\usepackage[numbers,sort&compress]{natbib}
\usepackage{fancyhdr}
\pagestyle{fancy}
\fancyhf{}
\fancyhead[L]{\small\itshape The Icosian Triad on $V_{600}$}
\fancyhead[R]{\small\thepage}
\fancyfoot[C]{\footnotesize\itshape Pre-peer-review open research preprint --- not independently validated.}
\renewcommand{\headrulewidth}{0.4pt}
\renewcommand{\footrulewidth}{0.4pt}

\newtheorem{theorem}{Theorem}
\newtheorem{proposition}[theorem]{Proposition}
\newtheorem{lemma}[theorem]{Lemma}
\newtheorem{corollary}[theorem]{Corollary}
\theoremstyle{definition}
\newtheorem{definition}[theorem]{Definition}
\theoremstyle{remark}
\newtheorem{remark}[theorem]{Remark}

\newcommand{\Z}{\mathbb{Z}}
\newcommand{\Q}{\mathbb{Q}}
\newcommand{\R}{\mathbb{R}}
\newcommand{\C}{\mathbb{C}}
\newcommand{\N}{\mathbb{N}}
\newcommand{\II}{\mathcal{I}}
\newcommand{\IIlip}{\mathcal{I}_{\mathrm{Lip}}}
\newcommand{\IImax}{\mathcal{I}_{\mathrm{max}}}
\newcommand{\GG}{\mathcal{G}}
\newcommand{\CC}{\mathcal{C}}
\newcommand{\NH}{\mathrm{N}_{\mathrm{H}}}
\newcommand{\NQ}{\mathrm{N}_{\mathbb{Q}(\sqrt{5})/\mathbb{Q}}}
\newcommand{\Zphi}{\Z[\varphi]}
\newcommand{\WHfour}{W(H_4)}

\title{\textbf{The Icosian Triad on $V_{600}$}\\[2pt]
\large A geometry-first derivation of $(\II, \GG, \CC)$ with $\sigma$-equivariant Galois structure\\[2pt]
\normalsize and Eichler--Hecke identification of the icosian theta L-function\\[6pt]
\small (Open research preprint --- pre-peer-review draft)}

\author{%
  Lee Smart\\[4pt]
  \textit{Institute of Vibrational Field Dynamics}\\[2pt]
  \texttt{contact@vibrationalfielddynamics.org}\\[2pt]
  \texttt{@vfd\_org}%
}
\date{May 2026}

\begin{document}
\maketitle
\thispagestyle{empty}

\begin{abstract}
\noindent
We give a geometry-first verification of three structural operators
attached to the vertex set $V_{600}$ of the regular $600$-cell, viewed
as the unit group of the maximal icosian order $\IImax$ over
$\Q(\sqrt{5})$ (Definition~\ref{def:order}). All universal arithmetic
claims (representation counts $r(\pi)$, multiplicativity, L-function
identity) are stated and proved for $\IImax$ and imported from the
classical Eichler--Hijikata / Eichler--Brandt theorems. The finite
explicit witnesses produced by our simulations live in the smaller
Lipschitz suborder $\IIlip = \Zphi^4 \subset \IImax$.

\smallskip\noindent\textbf{The triad.}
\emph{(i)~Substrate.} $V_{600}$ carries a multiplicity-free $9$-class symmetric
Bose--Mesner association scheme whose edge adjacency $A_1$ has nine distinct
eigenvalues
\[
\bigl\{\, 12,\;\; 6\varphi,\;\; 4\varphi,\;\; 3,\;\; 0,\;\; -2,\;\; 4-4\varphi,\;\; -3,\;\; 6-6\varphi \,\bigr\}\subset \Zphi
\]
with multiplicities $(1,4,9,16,25,36,9,16,4)$. The scheme is
generated by $A_1$ as an algebra over $\Q$.
\emph{(ii)~Closure.} The operator $\CC_\varphi := L_{V_{600}} + \varphi^{-2}\!\cdot\! I = (12+\varphi^{-2})\!\cdot\! I - A_1$
acts on each eigenspace $E_j$ as multiplication by the scalar $(12+\varphi^{-2}) - \lambda_j$.
\emph{(iii)~Generation.} The icosian quaternion norm
$\NH : \II \to \Zphi$ defined by $\NH(q) = q\bar q = q_0^2 + q_1^2 + q_2^2 + q_3^2$
is surjective onto the set of totally positive $\Zphi$-primes (up to units),
with explicit representation counts
\[
  r(\pi) \;:=\; \#\{\, q \in \IImax : \NH(q) = \pi\,\} \;=\; 120\bigl(1 + \NQ(\pi)\bigr)
  \qquad \text{for every }\Zphi\text{-prime }\pi,
\]
\emph{including} the inert prime $2$, where $r(2) = 600 = 120\cdot 5$ (no
anomaly: the maximal order is locally standard at $2$). The normalised count
$\widetilde r := r/120 = 1 + \NQ(\pi) = \sigma_1^{\Zphi}((\pi))$ is the
Eisenstein divisor coefficient. The bounded-range simulations of
\S\ref{sec:repro} are performed in the Lipschitz suborder
$\IIlip = \Zphi^4$ ($8$ units), which agrees with $\IImax$ away from $2$ but
is non-maximal at $2$, giving the suborder value $r_{\IIlip}(2)=24$ --- an
artifact of $\IIlip$, not a property of the icosian ring.

\smallskip\noindent\textbf{Galois equivariance.}
The Galois involution $\sigma : \varphi \mapsto 1-\varphi$ acts on $\Zphi$, on
the closure operator family $(\CC_\varphi, \CC_{1-\varphi})$, and on the
$A_1$-spectrum by pairing the eigenspaces $(E_1, E_8)$ and $(E_2, E_6)$ while
fixing $(E_0, E_3, E_4, E_5, E_7)$. The triad is $\sigma$-equivariant:
$\sigma \circ \NH = \NH \circ \widehat{\sigma}$, where $\widehat{\sigma}$ acts
coordinate-wise on icosian quaternions. The total dimension of the
$\sigma$-fixed eigenspace family equals
$\dim E_0 + \dim E_3 + \dim E_4 + \dim E_5 + \dim E_7 = 1+16+25+36+16 = 94$.

\smallskip\noindent\textbf{The L-function identity.}
The normalised representation function $\widetilde r = r/120$ is multiplicative:
$\widetilde r(\alpha\beta) = \widetilde r(\alpha)\,\widetilde r(\beta)$ for coprime
totally positive $\alpha, \beta \in \Zphi$. Since $\widetilde r(\pi) =
\sigma_1^{\Zphi}((\pi))$ at \emph{every} prime (including $2$), the icosian theta
L-function is the divisor-sum Dirichlet series, which factors \emph{exactly}:
\[
  L(\Theta_\II, s) \;:=\; \sum_{0 \neq \alpha \in \Zphi/\Zphi^\times}
  \widetilde r(\alpha)\!\cdot\! \NQ(\alpha)^{-s}
  \;=\; \zeta_K(s)\!\cdot\! \zeta_K(s-1)
\]
for $K = \Q(\sqrt{5})$ --- with \emph{no} local-$2$ correction
(Theorem~\ref{thm:euler-bridge}). Via the classical Dedekind factorisation
$\zeta_K(s) = \zeta(s)\!\cdot\! L(s,\chi_5)$, the Riemann zeta function
$\zeta(s)$ appears as one Dedekind factor of $L(\Theta_\II, s)$.

\smallskip\noindent\textbf{Status.}
The finite computational claims (spectrum, intersection-array, $9$-class
scheme axioms, finite witnesses, finite multiplicativity tests) are
verified by exact-arithmetic or controlled floating-point simulation
(see~\S\ref{sec:repro}). The maximal-order representation counts
$r(\pi)=120(1+N_Q(\pi))$ at every prime type (split, inert---including
$2$---and ramified) and the resulting exact identity
$L(\Theta_\II,s)=\zeta_K(s)\zeta_K(s-1)$ are verified by complete
Fincke--Pohst short-vector enumeration in the maximal icosian ring
(\texttt{sim\_theta\_maximal\_order.py}); they are also the classical
Eichler--Hijikata / Eichler--Brandt statements for a maximal order in a
totally definite quaternion algebra of class number one. We do not address
the Riemann Hypothesis or its Hilbert-modular analogue.
\end{abstract}

\bigskip

\noindent\fbox{%
\parbox{0.96\linewidth}{%
\textbf{Plain-language summary.}

\smallskip\noindent
The $600$-cell is a regular four-dimensional polytope with $120$ vertices,
each realised as a quaternion. These $120$ quaternions form a group: the
binary icosahedral group $2I$. The set is closed under a specific
non-commutative multiplication, and is contained in a quaternion algebra
defined over the field $\Q(\sqrt{5})$.

\smallskip\noindent
This paper derives, by direct computation and exact arithmetic, three
operators attached to $V_{600}$ — the \emph{substrate} ($V_{600}$ itself with
its full distance structure), the \emph{closure operator}
$\CC_\varphi$ (a specific shift of the graph Laplacian), and the
\emph{generation operator} $\NH$ (the quaternion norm map). Together these
three form a coherent \emph{triad} satisfying a Galois equivariance condition.

\smallskip\noindent
The main quantitative result is a closed-form count $r(\pi)$ of how many
icosian quaternions have norm equal to a given prime element $\pi$. From this
we identify the icosian theta function with an explicit weight-$2$ Hilbert
modular Eisenstein series over $\Q(\sqrt{5})$, whose L-function factors as
$\zeta_K(s)\cdot\zeta_K(s-1)$. The classical Riemann zeta function appears as
one Euler factor.

\smallskip\noindent
The paper does not prove the Riemann Hypothesis, the generalised Riemann
Hypothesis for $L(s,\chi_5)$, or any Hilbert-modular analogue.
}}

\bigskip

\section{Introduction}\label{sec:intro}

The vertex set $V_{600}$ of the regular $600$-cell has been studied from many
perspectives: as the binary icosahedral group $2I$, as a distance-transitive
graph, as the discrete unit group of a quaternion order, and as a substrate
for spectral analyses in the Vibrational Field Dynamics programme.
This paper gives a single self-contained derivation of three structural
operators on $V_{600}$ — substrate, closure, generation — and shows that they
form a Galois-equivariant triad whose theta L-function is an explicit Hilbert
modular Eisenstein series for $\Q(\sqrt{5})$.

The derivation is \emph{geometry-first}: the finite $V_{600}$ claims are
established by focused exact-arithmetic or controlled numerical
simulations. The universal arithmetic claims are imported from
classical Eichler--Hijikata / Eichler--Brandt theory and checked in
bounded finite ranges by the accompanying simulations. There are no
conditional hypotheses in the finite $V_{600}$ verification layer. The
trade-off is scope: we do not prove the Riemann Hypothesis or any of its
generalisations.

\subsection*{The triad}

The three operators we attach to $V_{600}$ are:

\begin{itemize}
\item \textbf{Substrate} $\II$: the icosian ring (a maximal $\Zphi$-order in
the quaternion algebra over $\Q(\sqrt{5})$), with $V_{600}$ as its group of
units. As a metric set, $V_{600}$ carries an Euclidean distance structure
giving rise to a $9$-class association scheme.

\item \textbf{Closure} $\CC_\varphi$: the operator $L_{V_{600}} + \varphi^{-2}\cdot I$
on the function space $\C[V_{600}]$. Here $L_{V_{600}} = 12\!\cdot\! I - A_1$ is
the graph Laplacian and $\varphi = (1+\sqrt 5)/2$. The constant $\varphi^{-2}$ is
fixed by the natural structure of $\Q(\sqrt{5})$ and the involution we will
identify below.

\item \textbf{Generation} $\GG = \NH$: the quaternion norm map
$\NH : \II \to \Zphi$ given by $\NH(q) = q\bar q = q_0^2 + q_1^2 + q_2^2 + q_3^2$.
Its image, restricted to totally positive elements, is exactly the set of
totally positive elements of $\Zphi$ (i.e., the universal representation
theorem for the icosian quadratic form).
\end{itemize}

The triad is glued together by Galois. The non-trivial automorphism
$\sigma : \Q(\sqrt 5) \to \Q(\sqrt 5)$ takes $\varphi$ to $1-\varphi$ and
extends to a coordinate-wise involution $\widehat\sigma$ on
$\II = (\Zphi)^4$-quaternions. We will show that $\sigma$ commutes with
$\NH$ ($\sigma\NH = \NH\widehat\sigma$), and that $\sigma$ permutes the
$A_1$-eigenspaces of $V_{600}$ by pairing eigenvalues with their Galois
conjugates.

\subsection*{The L-function identification}

The icosian theta series counts representations of $\Zphi$-elements as sums of
four icosian quaternion squares. Sims~11--13 establish (with $\widetilde r := r/120$,
the count normalised by the $120$ units of the maximal order):

\begin{itemize}
  \item $\widetilde r(\pi) = 1 + \NQ(\pi) = \sigma_1^{\Zphi}((\pi))$ for
    \emph{every} totally positive $\Zphi$-prime $\pi$ --- split, inert
    (including $2$, where $\widetilde r(2)=5$), and ramified
    (Eichler--Hecke closed form for the maximal order).
  \item $\widetilde r$ is multiplicative on coprime totally positive
    $\Zphi$-elements: $\widetilde r(\alpha\beta) = \widetilde r(\alpha)\widetilde r(\beta)$.
  \item The Dirichlet series of $\widetilde r$ factors \emph{exactly} as
    $\zeta_K(s)\!\cdot\!\zeta_K(s-1)$, $K = \Q(\sqrt 5)$ --- no local-$2$ correction.
\end{itemize}

The classical Dedekind factorisation $\zeta_K(s) = \zeta(s)\!\cdot\! L(s,\chi_5)$
then identifies the Riemann zeta function $\zeta(s)$ as one Euler factor of
the icosian theta L-function:
\[
  L(\Theta_\II, s) \;=\; \zeta(s)\!\cdot\! L(s,\chi_5)\!\cdot\! \zeta(s-1)\!\cdot\! L(s-1, \chi_5)
  \qquad \text{(exactly, for the maximal order).}
\]
This is the geometry-first realisation of the cascade Euler bridge
between the classical Riemann zeta function and the Hilbert modular
content of $\Q(\sqrt 5)$ via the icosian quaternion order.

\subsection*{What this paper does NOT do}

It does not prove the Riemann Hypothesis, the generalised RH for $L(s,\chi_5)$,
or the Hilbert-modular RH analogue for $\Q(\sqrt 5)$. The L-function we
identify has the same zeros as $\zeta_K(s)\cdot\zeta_K(s-1)$, which lie on
$\mathrm{Re}(s)=1/2$ and $\mathrm{Re}(s)=3/2$ if and only if classical RH and
GRH hold. We do not address those questions; they remain open and are
classical.

\subsection*{Organisation}

Sections~\ref{sec:substrate} through~\ref{sec:generation} introduce the
substrate, closure, and generation operators. Section~\ref{sec:sigma}
establishes Galois equivariance and identifies the $94$ $\sigma$-fixed
dimensions. Section~\ref{sec:reps} gives the closed-form representation
counts and multiplicativity. Section~\ref{sec:Lfunction} proves the L-function
identification. Section~\ref{sec:open} states open questions and the precise
boundary between what we derive and what remains classical. Reproduction
information is in Section~\ref{sec:repro}.

\subsection*{Reproduction}

Every finite computational claim in this paper is backed by a focused
simulation. Universal arithmetic claims are supported by the cited
classical theorems (Eichler--Hijikata, Eichler--Brandt) and checked
against the finite simulation outputs where applicable. The complete
reproduction infrastructure is in the accompanying repository
\texttt{repro/}, with thirteen \texttt{sim\_*.py} scripts numbered
1--13. The acceptance criteria for each sim are documented in
\texttt{repro/sims/FINDINGS.md}. All numerical claims in the paper are
verbatim from \texttt{repro/output/*.json}.

\bigskip

\section{Substrate: $V_{600}$ and its association scheme}\label{sec:substrate}

\subsection{Construction of $V_{600}$}

We work in the quaternion algebra $\R \cdot 1 \oplus \R\!\cdot\! i \oplus \R\!\cdot\! j \oplus \R\!\cdot\! k$
with $i^2=j^2=k^2=ijk=-1$. The unit icosians form the $120$-element set
\begin{align}
V_{600} \;=\;
\;& \bigl\{ \pm 1,\pm i,\pm j,\pm k\bigr\} && \text{(8 axis vertices)} \notag \\
+ \;& \bigl\{ (\pm 1 \pm i \pm j \pm k)/2 \bigr\} && \text{(16 half-cube vertices)} \notag \\
+ \;& \bigl\{ \text{even permutations of } (\pm \tfrac{\varphi}{2},\pm \tfrac{1}{2},\pm \tfrac{1}{2\varphi}, 0) \bigr\} && \text{(96 snub vertices)} \notag
\end{align}
where $\varphi = (1+\sqrt 5)/2$ is the golden ratio. The set is closed under
quaternion multiplication and inversion; it is the binary icosahedral group
$2I$. Each vertex has $\|v\|^2 = 1$.

\subsection{The 9-class Euclidean association scheme}

The pairwise squared Euclidean distances on $V_{600}$ take exactly nine
distinct values
\[
  d^2 \in \bigl\{0,\; 1/\varphi^2,\; 1,\; 1+1/\varphi^2,\; 2,\; 2+1/\varphi^2,\; 3,\; 3+1/\varphi^2,\; 4 \bigr\}.
\]
For $i=0,\dots,8$, let $A_i \in M_{120}(\Z)$ denote the $0/1$-incidence
operator at distance $d_i$. $A_0 = I$, $\sum_i A_i = J$ (all-ones matrix),
and the $A_i$ form a commutative algebra under matrix multiplication.

\begin{theorem}[Association scheme; sim 4]\label{thm:scheme}
$\{A_0,\dots,A_8\}$ spans a $9$-dimensional symmetric Bose--Mesner algebra:
\begin{enumerate}
  \item $A_i A_j = \sum_k p_{ij}^k A_k$ with non-negative integer
    intersection numbers $p_{ij}^k$.
  \item The valency vector is $(k_0,\dots,k_8) = (1,12,20,12,30,12,20,12,1)$.
  \item Let $\{E_j\}_{j=0}^8$ be the system of primitive orthogonal
    idempotents (projectors onto the eigenspaces of $A_1$). Then
    $A_i = \sum_j P_{ij} E_j$ and $E_j = (1/120) \sum_i Q_{ij} A_i$ with
    $9\times 9$ matrices $P, Q$ satisfying $PQ = QP = 120\cdot I$ exactly.
  \item The Krein parameters $q_{ij}^k$ are all non-negative.
\end{enumerate}
The scheme has multiplicity-free permutation representation.
\end{theorem}

\begin{proof}
Direct computation in \texttt{sim\_association\_scheme.py}; fifteen acceptance
checks pass with exact arithmetic in $\Q(\sqrt 5)$ or with floating-point
tolerance $10^{-9}$.
\end{proof}

\subsection{The spectrum of $A_1$}

\begin{theorem}[Spectrum; sims 1--2, 5]\label{thm:spectrum}
The edge adjacency $A_1$ has exactly $9$ distinct eigenvalues lying in
$\Zphi$. With the multiplicity sequence
$(\dim E_j) = (1,4,9,16,25,36,9,16,4)$ summing to $120$, the spectrum is
\[
\begin{array}{|c|c|c|c|}
\hline
j & \lambda_j \in \Zphi & \dim E_j & \text{numerical value} \\
\hline
0 & 12          & 1   & +12 \\
1 & 6\varphi    & 4   & +9.7082 \\
2 & 4\varphi    & 9   & +6.4721 \\
3 & 3           & 16  & +3 \\
4 & 0           & 25  & \phantom{+}0 \\
5 & -2          & 36  & -2 \\
6 & 4-4\varphi  & 9   & -2.4721 \\
7 & -3          & 16  & -3 \\
8 & 6-6\varphi  & 4   & -3.7082 \\
\hline
\end{array}
\]
$A_1$ generates the Bose--Mesner algebra as an algebra over $\Q$:
each $A_i$ ($i = 0, \dots, 8$) is expressible as a polynomial in $A_1$.
\end{theorem}

\begin{proof}
\texttt{sim\_geometry\_first.py} verifies the multiplicities by computing
the nine $\WHfour$-equivariant fingerprints under a set of five generators
of $\WHfour$. \texttt{sim\_p\_polynomial\_structure.py} explicitly
constructs polynomials $p_i$ with $A_i = p_i(A_1)$ to floating-point
precision $\sim 10^{-10}$. We do not claim the standard BCN \emph{narrow}
$P$-polynomial property, which would require a three-term tridiagonal
recurrence; that recurrence fails for this scheme
(Remark~\ref{rem:not-narrow-DRG}).
\end{proof}

\begin{remark}\label{rem:not-narrow-DRG}
The $9$-class Euclidean scheme is generated as an algebra by $A_1$ but is
\emph{not} a narrow distance-regular graph in the sense of
Brouwer--Cohen--Neumaier: the intersection numbers do not satisfy
$a_i + b_i + c_i = k_1 = 12$ for interior shells, because $A_1 \cdot A_s$
has off-band contributions in this scheme.
Independently, the graph-distance partition of $V_{600}$'s $1$-skeleton has
only six classes (valencies $1,12,32,42,32,1$); that partition is also not a
strict DRG because the antipode's neighbours split unevenly among the
graph-distance-$4$ shell. See \texttt{sim\_graph\_distance\_scheme.py}.
\end{remark}

\bigskip

\section{Closure operator $\CC_\varphi$}\label{sec:closure}

\begin{definition}\label{def:closure}
The \emph{closure operator} on $V_{600}$ is
\[
  \CC_\varphi \;:=\; L_{V_{600}} + \varphi^{-2}\!\cdot\! I
  \;=\; (12 + \varphi^{-2})\!\cdot\! I - A_1,
\]
where $L_{V_{600}} = 12\!\cdot\! I - A_1$ is the (combinatorial) graph
Laplacian of the $A_1$-graph.
\end{definition}

\begin{theorem}[Closure spectrum; sim 8]\label{thm:closure-spec}
$\CC_\varphi$ commutes with $A_1$. On each $A_1$-eigenspace $E_j$, $\CC_\varphi$
acts as multiplication by the scalar
\[
  \mu_j \;=\; (12 + \varphi^{-2}) - \lambda_j \;\in\; \Zphi[\varphi^{-2}]
       \;=\; \Zphi[\,\textstyle\frac{1}{1+\varphi}\,]
       \;\subset\; \Q(\sqrt 5).
\]
The Galois conjugate operator $\CC_{1-\varphi} := L_{V_{600}} + \varphi^2\!\cdot\! I$
acts on $E_j$ as $\mu_j^\sigma := (12 + \varphi^2) - \lambda_j$. These are
distinct operators on $\C[V_{600}]$ with distinct spectra; $(\CC_\varphi, \CC_{1-\varphi})$
forms a $\sigma$-pair of operators.
\end{theorem}

\begin{proof}
$\CC_\varphi = (12 + \varphi^{-2}) I - A_1$ is a polynomial in $A_1$, hence
shares eigenspaces with $A_1$ and has eigenvalues $\mu_j$. Numerical
verification in \texttt{sim\_closure\_sigma\_action.py} shows all nine
predicted scalars match measurement to machine precision.
\end{proof}

\bigskip

\section{Generation operator $\GG = \NH$}\label{sec:generation}

\subsection{The icosian quaternion norm}

\begin{definition}[icosian order model]\label{def:order}
This paper uses two related orders inside the quaternion algebra over
$\Q(\sqrt 5)$:
\begin{itemize}
\item the \emph{Lipschitz $\Zphi$-order}
  $\IIlip := \Zphi \oplus \Zphi\!\cdot\! i \oplus \Zphi\!\cdot\! j \oplus \Zphi\!\cdot\! k$,
  a rank-$4$ $\Zphi$-module with integer-coordinate quaternions only;
\item the \emph{maximal icosian order} $\IImax$
  \cite{conwaySloane}, a strict superset of $\IIlip$ that includes
  half-integer combinations (specifically, the icosian Hurwitz-style
  elements such as $(1+i+j+k)/2$ and the half-integer snub icosians).
\end{itemize}

The vertex set $V_{600}$ lies in $\IImax$ as its unit group of order
$120$. Of these, only the $8$ axis vertices $(\pm 1, 0, 0, 0)$ and
permutations and the $16$ half-cube vertices $(\pm 1/2, \pm 1/2, \pm 1/2, \pm 1/2)$
sit in $\IIlip$ (after rescaling to clear the half-integers from
$\IIlip$ they live in $\IIlip$; the $96$ snub vertices do not).

\textbf{Scope convention.} Throughout this paper:
\begin{itemize}
\item All \emph{universal} arithmetic claims (representation count
  $r(\pi) = 120(1 + N_Q(\pi))$ at every prime; multiplicativity of
  $\widetilde r = r/120$ on coprime ideals; the exact L-function identity)
  are stated for $\IImax$, verified by complete short-vector enumeration in
  $\IImax$ (\texttt{sim\_theta\_maximal\_order.py}), and agree with the
  classical Eichler--Hijikata / Eichler--Brandt theorems
  \cite{eichler1958, hijikata1974}.
\item All \emph{finite} explicit witness searches (Theorem~\ref{thm:univ-rep},
  sim~9) are conducted in the Lipschitz suborder $\IIlip$ and are
  reported as bounded reproducibility checks for the classical
  universal statement.
\end{itemize}
\end{definition}

\begin{remark}[Theorem imports vs finite simulations]
The reader should keep two complementary streams in mind:
\begin{enumerate}
\item \textbf{Finite, exact-arithmetic claims} (the $9$-class scheme,
  the spectrum in $\Zphi$, the $94$-dimensional $\sigma$-fixed sum,
  the $20$ bounded witnesses for ideals of small norm, the six
  multiplicativity tests, the $s$-value L-function tests at
  $s \in \{2.5, 3, 3.5, 4, 5\}$) are verified by the simulations in
  \texttt{repro/sims/} to either exact $\Q(\sqrt 5)$ rational arithmetic
  or controlled floating-point tolerance $\sim 10^{-10}$.
\item \textbf{Universal arithmetic claims} (representation counts and
  multiplicativity on all coprime $\IImax$-ideals; prime-power Euler
  factors; the complete L-function identity) are imports from classical
  Eichler--Hijikata / Eichler--Brandt theory applied to the icosian
  quaternion order. Our finite simulations serve as reproducibility
  checks and witnesses for these classical theorems, not as their
  proofs.
\end{enumerate}
\end{remark}

\begin{definition}\label{def:NH}
For $q = q_0 + q_1 i + q_2 j + q_3 k \in \IImax$ (or in the sublattice
$\IIlip$ if all $q_\nu \in \Zphi$), the quaternion norm is
\[
  \NH(q) \;:=\; q\bar q \;=\; q_0^2 + q_1^2 + q_2^2 + q_3^2 \;\in\; \Zphi.
\]
$\NH$ is multiplicative and takes totally positive values.
\end{definition}

\begin{theorem}[Bounded explicit witnesses; sim 9]\label{thm:univ-rep}
For every totally positive $\Zphi$-prime $\pi$ whose underlying rational
prime $p$ lies in $\{2,3,5,7,11,13,17,19,23,29,31,37,41,43,47\}$, there
exists $q \in \IIlip$ (i.e., all four coordinates in $\Zphi$ with
$|a|, |b| \leq 4$) such that $\NH(q) = \pi$. Explicit witnesses are
given for the $20$ canonical totally positive prime representatives
above these rational primes.

The universal statement (existence of $q \in \IImax$ for every
totally positive $\Zphi$-prime) is the classical universal-representation
theorem for the icosian quadratic form on the maximal icosian order,
which follows from the Hasse--Minkowski local-global principle together
with the local-density computations of Eichler~\cite{eichler1958} and
Hijikata~\cite{hijikata1974}. Our explicit witnesses verify this
classical result by producing representatives that lie in the smaller
Lipschitz sublattice $\IIlip \subset \IImax$ for the bounded range
stated.
\end{theorem}

\begin{proof}
\texttt{sim\_generation\_operator.py} performs an exhaustive meet-in-the-middle
search over the Lipschitz lattice with $|a|,|b| \leq 4$ per coordinate and
finds a witness $q$ for each of the $20$ canonical primes. Verification is
exact.
\end{proof}

The universality theorem for the icosian quadratic form over $\Zphi$ is
classical (Eichler--Hijikata). We re-derive it here at the level of explicit
witnesses, as needed for the Galois-equivariance computation.

\bigskip

\section{Galois equivariance: the triad isomorphism}\label{sec:sigma}

\subsection{The action of $\sigma$ on $\II$ and $V_{600}$}

Let $\sigma : \Q(\sqrt 5) \to \Q(\sqrt 5)$ be the non-trivial automorphism
$\varphi \mapsto 1-\varphi$. It extends to a coordinate-wise involution
\[
  \widehat\sigma : \II \to \II, \qquad
  \widehat\sigma(q_0 + q_1 i + q_2 j + q_3 k) \;=\; \sigma(q_0) + \sigma(q_1) i + \sigma(q_2) j + \sigma(q_3) k.
\]
$\widehat\sigma$ is an involution and a $\Z$-linear map; it commutes with
quaternion multiplication.

\begin{remark}
The Wikipedia embedding of $V_{600}$ as the $120$-vertex set above is
\emph{not} closed under $\widehat\sigma$ as a set of points in $\R^4$: only
the $8$ axis and $16$ half-cube vertices are coordinate-wise $\sigma$-fixed.
Each snub vertex maps under $\widehat\sigma$ to a snub vertex in the
\emph{odd}-permutation class. The icosian ring $\II$ itself is closed under
$\widehat\sigma$, but the geometric realisation chosen above places $V_{600}$
and $\widehat\sigma(V_{600})$ on different sides of an $\widehat\sigma$-orbit
in $\II$. We therefore work spectrally, where the $\sigma$-pairing structure
is canonical and embedding-independent.
\end{remark}

\subsection{$\sigma$-pairing of the $A_1$-spectrum}

\begin{theorem}[$\sigma$-pairing; sim 8]\label{thm:sigma-pair}
With the spectrum and multiplicities of Theorem~\ref{thm:spectrum}:
\begin{itemize}
  \item \textbf{$\sigma$-fixed eigenvalues:} $\{12, 3, 0, -2, -3\} \subset \Z$ are
    fixed by $\sigma$; the eigenspaces $E_0, E_3, E_4, E_5, E_7$ are
    individually $\sigma$-invariant (as $\Zphi$-modules of A1-eigenvectors with
    eigenvalue fixed by $\sigma$).
  \item \textbf{$\sigma$-paired eigenvalues:}
    $\sigma(6\varphi) = 6-6\varphi$, $\sigma(4\varphi) = 4-4\varphi$. The
    eigenspaces $E_1 \leftrightarrow E_8$ (each of dimension $4$) and
    $E_2 \leftrightarrow E_6$ (each of dimension $9$) form Galois-conjugate
    pairs.
\end{itemize}
The total dimension of the $\sigma$-fixed eigenspaces is
\[
  \dim E_0 + \dim E_3 + \dim E_4 + \dim E_5 + \dim E_7
  \;=\; 1 + 16 + 25 + 36 + 16 \;=\; 94,
\]
and the total dimension of one side of the $\sigma$-paired class is
$\dim E_1 + \dim E_2 = 4 + 9 = 13$. The full decomposition is $94 + 2\cdot 13 = 120$.
\end{theorem}

\begin{proof}
The $A_1$-eigenvalues are in $\Zphi$ and the characteristic polynomial of
$A_1$ has integer coefficients; Galois automorphisms therefore permute the
roots while preserving multiplicities. Direct computation in
\texttt{sim\_closure\_sigma\_action.py} verifies the pairing pattern and
dimension count.
\end{proof}

\subsection{Antipodal pairing of $(E_3, E_7)$}

The dimension-$16$ eigenspaces $E_3$ and $E_7$ have $\sigma$-fixed eigenvalues
$\pm 3$ and so are not paired by $\sigma$. They are nevertheless paired by a
different natural operator: the antipodal incidence $A_8$ (the operator
selecting the unique vertex at maximal Euclidean distance).

\begin{proposition}[Antipodal pairing; sim 8]\label{prop:antipodal}
$A_8$ acts as $-I$ on $E_3$ and as $+I$ on $E_7$. The pairing
$(E_3, E_7)$ is an antipodal $\Z_2$-grading, distinct from the
Galois $\sigma$-pairings of $(E_1, E_8)$ and $(E_2, E_6)$.
\end{proposition}

\subsection{$\sigma$-equivariance of $\NH$}

\begin{theorem}[Triad isomorphism; sim 10]\label{thm:triad-iso}
For every $q \in \IImax$ (and so in particular for every $q \in \IIlip$),
\[
  \NH\bigl(\widehat\sigma(q)\bigr) \;=\; \sigma\bigl(\NH(q)\bigr).
\]
Consequently:
\begin{itemize}
  \item \textbf{Inert primes} ($p \equiv \pm 2 \bmod 5$) admit
    $\widehat\sigma$-fixed icosian witnesses with coordinates in $\Z$
    (Lagrange-style four-squares representation in $\Z$).
  \item \textbf{Split prime pairs} $(\pi, \sigma(\pi))$ above rational
    primes $p \equiv \pm 1 \bmod 5$ correspond to $\widehat\sigma$-conjugate
    icosian pairs $(q, \widehat\sigma(q))$ with $\NH(q) = \pi$ and
    $\NH(\widehat\sigma(q)) = \sigma(\pi)$.
  \item \textbf{Ramified prime} above $5$: the generator $2 + \varphi$ of
    the ramified ideal has $\NH$-witness in $\II$ with $\sigma$-twisted
    structure.
\end{itemize}
\end{theorem}

\begin{proof}
$\NH(q) = q_0^2 + q_1^2 + q_2^2 + q_3^2$ is a polynomial in the coordinates
of $q$ with integer coefficients; therefore $\NH$ commutes with any
coordinate-wise field automorphism. Verification on the twenty witnesses of
Theorem~\ref{thm:univ-rep} in \texttt{sim\_triad\_isomorphism.py} confirms
the identity exactly. The classification of inert, split, ramified primes
of $\Zphi$ is the classical Dedekind theorem.
\end{proof}

\bigskip

\section{Representation numbers $r(\pi)$}\label{sec:reps}

\subsection{Closed form for primes}

\begin{theorem}[Closed-form representation count; Eichler--Hijikata + sim 13]\label{thm:rpi}
For \emph{every} totally positive $\Zphi$-prime $\pi$ --- split, inert, or
ramified ---
\[
  r(\pi) \;:=\; \#\{\,q \in \IImax : \NH(q) = \pi\,\}
  \;=\; 120\!\cdot\!(1 + \NQ(\pi)),
\]
in the maximal icosian order $\IImax$. In particular at the inert prime
$\pi = 2 \in \Zphi$ ($\NQ(2)=4$) the count is $r(2) = 600 = 120\cdot 5$; there
is \emph{no} local-$2$ anomaly. Equivalently $\widetilde r(\pi) := r(\pi)/120
= 1 + \NQ(\pi) = \sigma_1^{\Zphi}((\pi))$ for every prime.
\end{theorem}

\begin{proof}
The closed form is the Eichler--Hijikata representation-number formula for a
maximal order in a totally definite quaternion algebra over $K$ of class
number one~\cite{eichler1958,hijikata1974}: the local density at each finite
prime $\mathfrak p$ (split, inert, or ramified) is the standard
$1 + \NQ(\mathfrak p)$, since the algebra $B=(-1,-1/K)$ is unramified at every
finite prime, and the global constant is the order's unit-group size $120$
(the $120$ unit icosians). Our \texttt{sim\_theta\_maximal\_order.py} verifies
this by complete Fincke--Pohst short-vector enumeration in $\IImax$: split
primes ($N=11,19,29$), the inert primes $2$ ($r=600$) and $3$ ($r=1200$), and
the ramified prime above $5$ ($r=720$) all give $r/120 = 1+\NQ$, and the inert-$2$
\emph{powers} give $r((2)^k)/120 = 1,5,21,85 = \sigma_1((2)^k)$ --- the standard
Eisenstein local factor, so no correction.

\emph{Remark (the earlier ``$r(2)=24$'').} The bounded-range enumerator
\texttt{sim\_representation\_numbers.py} counts the \emph{Lipschitz suborder}
$\IIlip = \Zphi^4$ (only $8$ units), which agrees with $\IImax$ at odd primes
(giving $r_{\IIlip}=8(1+\NQ)$, i.e.\ the same normalised $\widetilde r$) but is
non-maximal at $2$, where it returns $r_{\IIlip}(2)=24$. That value, and the
local-$2$ correction it would induce, are artifacts of the suborder and do not
hold for the icosian ring.
\end{proof}

\subsection{Multiplicativity}

\begin{theorem}[Multiplicativity; Eichler--Brandt + sim 12]\label{thm:mult}
For coprime totally positive $\alpha, \beta \in \Zphi$, the normalised count
$\widetilde r := r/120$ (maximal order) is multiplicative:
\[
  \widetilde r(\alpha\!\cdot\!\beta) \;=\; \widetilde r(\alpha)\!\cdot\! \widetilde r(\beta),
  \qquad\text{equivalently}\qquad
  r(\alpha\beta) \;=\; \tfrac{1}{120}\,r(\alpha)\,r(\beta).
\]
Since the maximal order is locally standard at every prime (including $2$),
no prime is excluded.
\end{theorem}

\begin{proof}
This is a special case of the Eichler--Brandt
theorem~\cite{eichler1958,hijikata1974} applied to the icosian
quaternion order's theta series, which is multiplicative on coprime
$\Zphi$-ideals coprime to the discriminant. Our
\texttt{sim\_multiplicativity.py} verifies the identity explicitly for
six coprime pairs $(\alpha, \beta)$ with composite norm up to $209$:
the inert--inert pair $(2,3)$, the inert--ramified $(2, 2+\varphi)$ and
$(3, 2+\varphi)$, the inert--split $(3, 3+\varphi)$, the split--split
$(3+\varphi, 4+\varphi)$ above distinct rational primes $11$ and $19$,
and the split--$\sigma$-conjugate pair $(3+\varphi, 4-\varphi)$ above
$11$. The all-pair statement is the cited Eichler--Brandt theorem.
\end{proof}

\bigskip

\section{The icosian theta L-function}\label{sec:Lfunction}

\subsection{Definition}

\begin{definition}\label{def:LTheta}
The icosian theta L-function is the Dirichlet series over $\Zphi$-ideals
\[
  L(\Theta_\II, s) \;:=\; \sum_{0 \neq \mathfrak{a}\, \subseteq\, \Zphi}
  \frac{\widetilde r(\mathfrak{a})}{\NQ(\mathfrak{a})^{s}},
\]
where the sum is over distinct nonzero integral $\Zphi$-ideals
$\mathfrak{a}$ and $\widetilde r(\mathfrak{a}) := r(\alpha) / 8$ for any
generator $\alpha$ of $\mathfrak{a}$.
\end{definition}

\begin{remark}[Well-definedness on ideals]\label{rem:well-defined}
$\Zphi$ is a principal ideal domain (class number $1$ for $K =
\Q(\sqrt 5)$), so every nonzero integral ideal $\mathfrak{a}$ is
principal: $\mathfrak{a} = (\alpha)$ for some $\alpha \in \Zphi$, unique
up to multiplication by a unit of $\Zphi$. The representation count
$r(\alpha) = \#\{q \in \IImax : \NH(q) = \alpha\}$ is invariant under
multiplication of $\alpha$ by a unit of $\IImax$ acting on $q$, hence
$r(u\alpha) = r(\alpha)$ for $u \in \Zphi^\times = \{\pm \varphi^n\}_{n\in\Z}$
(the totally-positive subset specifically; the multiplication-by-$-1$
case is absorbed into the count via the $\pm$ symmetry of the
quaternion norm form). The factor $1/8$ removes the action of the
size-$8$ subgroup
$\{\pm 1, \pm i, \pm j, \pm k\} \subset V_{600}$ of unit icosians
preserving any single representative direction. The resulting
$\widetilde r(\mathfrak{a})$ depends only on the ideal $\mathfrak{a}$,
not the chosen generator $\alpha$.
\end{remark}

\subsection{The Eichler--Hecke identity}

\begin{theorem}[L-function identification, Eichler--Brandt + sim 13]\label{thm:euler-bridge}
Let $K = \Q(\sqrt 5)$ and $\zeta_K(s) = \prod_{\mathfrak{p}} (1 - \NQ(\mathfrak{p})^{-s})^{-1}$
the Dedekind zeta function of $K$. For the \emph{maximal} icosian order $\IImax$,
\[
  L(\Theta_\II, s) \;=\; \zeta_K(s)\!\cdot\! \zeta_K(s-1)
  \qquad\text{(exactly; no local correction).}
\]
The local Euler factor at \emph{every} prime $\mathfrak{p}$ --- split, inert
(including $\mathfrak{p}_2=(2)$), and ramified --- coincides with that of
$\zeta_K(s)\zeta_K(s-1)$.
\end{theorem}

\begin{proof}
By Theorem~\ref{thm:rpi}, $\widetilde r(\pi)=1+\NQ(\pi)=\sigma_1^{\Zphi}((\pi))$
at every prime, and by Theorem~\ref{thm:mult} $\widetilde r$ is multiplicative;
hence $\widetilde r(\mathfrak a)=\sigma_1^{\Zphi}(\mathfrak a)$ for every
$\Zphi$-ideal $\mathfrak a$. The Dirichlet series identity
$\sum_{\mathfrak a}\sigma_1(\mathfrak a)\NQ(\mathfrak a)^{-s}
=\zeta_K(s)\zeta_K(s-1)$ is the standard Eisenstein divisor-sum identity for
$K$. The only delicate prime is the inert $\mathfrak p_2=(2)$, where the
algebra $B=(-1,-1/K)$ is nonetheless unramified; the maximal-order local factor
there is the standard $(1-4^{-s})^{-1}(1-4^{1-s})^{-1}$, confirmed by complete
Fincke--Pohst enumeration giving $\widetilde r((2)^k)=\sigma_1((2)^k)=1,5,21,85$
(\texttt{sim\_theta\_maximal\_order.py}). Hence no local-$2$ correction occurs.

\emph{Retraction.} Earlier drafts (and \texttt{sim\_theta\_L\_function.py})
asserted a correction factor $C_2(s)=1-2\cdot2^{-s}+2^{2-2s}$ from a value
$r(2)=24$. That value is the count in the \emph{Lipschitz suborder} $\IIlip$
(8 units), which is non-maximal at $2$; for the maximal icosian ring
$r(2)=600=120\cdot5$ and $C_2\equiv 1$. The earlier $C_2$ is withdrawn. (It was
also dimensionally ill-formed: an inert prime of norm $4$ admits only a local
factor in $4^{-s}$, not the $2^{-s}$ written there.)
\end{proof}

\subsection{Appearance of the Riemann zeta function}

\begin{corollary}[Euler bridge]\label{cor:euler-bridge}
Via the classical Dedekind factorisation $\zeta_K(s) = \zeta(s)\!\cdot\! L(s,\chi_5)$,
the icosian theta L-function admits the factorisation
\[
  L(\Theta_\II, s) \;=\; \zeta(s)\!\cdot\! L(s,\chi_5)\!\cdot\! \zeta(s-1)\!\cdot\! L(s-1,\chi_5),
\]
in which the Riemann zeta function $\zeta(s)$ appears as one Euler factor of
the geometry-derived icosian theta.
\end{corollary}

This is the bridge from the icosian quaternion structure on $V_{600}$ to
classical analytic number theory. The bridge is structural and Eisenstein-style;
it does not address the zero-distribution questions (RH, GRH).

\bigskip

\section{What this paper does not do}\label{sec:open}

\subsection{The Riemann Hypothesis and its generalisations}

Theorem~\ref{thm:euler-bridge} identifies $L(\Theta_\II, s)$ with
$\zeta_K(s)\!\cdot\!\zeta_K(s-1)$ exactly. The zero set of
$L(\Theta_\II, s)$ therefore coincides with the union of:
\begin{enumerate}
  \item Zeros of $\zeta(s)$ — the classical Riemann zeros;
  \item Zeros of $L(s, \chi_5)$ — the generalised Riemann zeros of $\chi_5$;
  \item Zeros of $\zeta(s-1)$ and $L(s-1, \chi_5)$ — shifts of the above.
\end{enumerate}
The classical Riemann Hypothesis says (1) lies on $\mathrm{Re}(s) = 1/2$. The
generalised Riemann Hypothesis for Dirichlet L-functions says (2) lies on
$\mathrm{Re}(s) = 1/2$. Both are open. Their truth would place every zero of
$L(\Theta_\II, s)$ on $\mathrm{Re}(s) = 1/2$ or $\mathrm{Re}(s) = 3/2$. The
present paper does not address these questions.

\subsection{The identification of $\tau_\mathrm{spec}$}

In~\cite{existenceClosure}, two involutions on $V_{600}$ are used: $\tau_\mathrm{ico}$
with $\dim \mathrm{Fix} = 94$ (the Galois involution of this paper) and
$\tau_\mathrm{spec}$ with $\dim \mathrm{Fix} = 116$. The present paper identifies
$\tau_\mathrm{ico}$ as the Galois $\sigma$-action of Theorem~\ref{thm:sigma-pair}.
We do not identify $\tau_\mathrm{spec}$ geometrically here; that is a separate
combinatorial question.

\subsection{Higher Hecke operators}

The Brandt matrices of the icosian order at higher prime powers act on
$\C[V_{600}]$ and admit further spectral decomposition. We do not develop the
Hecke-operator framework here.

\bigskip

\section{Appendix: the local factor at $2$ (maximal order; no correction)}\label{sec:local-2}

This appendix records the local Euler factor at the inert prime
$\mathfrak{p}_2=(2)$ for the \emph{maximal} icosian order, and explains why
an earlier draft's correction factor $C_2(s)$ was spurious.

The rational prime $2$ is inert in $\Zphi$: $(2)$ is prime of rational norm
$\NQ((2))=4$. The algebra $B=(-1,-1/K)$ is ramified only at the two
\emph{infinite} places, hence unramified at every finite prime, $2$ included.
For a maximal order in such an algebra the local theta factor at a prime of
norm $q$ is the standard
\[
  L_{\mathfrak p}(s) \;=\; (1 - q^{-s})^{-1}(1 - q^{1-s})^{-1},
\]
so at $\mathfrak p_2$ (with $q=4$) it is $(1-4^{-s})^{-1}(1-4^{1-s})^{-1}$,
predicting $r(2)=120\cdot(1+4)=600$ and, more fully,
$\widetilde r((2)^k)=\sigma_1((2)^k)=(4^{k+1}-1)/3 = 1,5,21,85,\dots$.
Complete Fincke--Pohst enumeration in the maximal icosian ring
(\texttt{sim\_theta\_maximal\_order.py}) confirms these exactly:
$r((2)^k)=120\cdot(1,5,21,85)=120,600,2520,10200$. The generating function
$\sum_k \sigma_1((2)^k)\,(4^{-s})^k = \big[(1-4^{-s})(1-4^{1-s})\big]^{-1}$ is
precisely the standard factor, so the local-$2$ correction is
\[
  C_2(s) \;=\; 1.
\]

\emph{What went wrong before.} A previous draft took $r(2)=24$ and built a
factor $C_2(s)=1-2\cdot2^{-s}+2^{2-2s}$. The value $24$ is the count in the
\emph{Lipschitz suborder} $\IIlip=\Zphi^4$ (which has only $8$ units and is
non-maximal at $2$), used as a bounded-range enumeration proxy; it is not the
maximal-order count. The stated factor was moreover ill-formed dimensionally:
an inert prime of norm $4$ admits a local factor only in $4^{-s}$, whereas
$1-2\cdot2^{-s}+\dots$ contains a bare $2^{-s}$. Both the value and the factor
are withdrawn; the maximal-order answer is the clean $C_2=1$ above.

\bigskip

\section{Reproduction}\label{sec:repro}

All numerical claims in this paper are produced by the simulations listed in
\texttt{repro/sims/}. The mapping from theorem to simulation is:

\begin{center}
\begin{tabular}{|c|l|l|}
\hline
\textbf{Theorem} & \textbf{Sim} & \textbf{Output} \\
\hline
\ref{thm:scheme} (Scheme axioms)        & \texttt{sim\_association\_scheme.py}    & \texttt{association\_scheme.json} \\
\ref{thm:spectrum} (Spectrum)            & \texttt{sim\_geometry\_first.py}        & \texttt{geom\_first\_results.json} \\
                                          & \texttt{sim\_irrep\_identification.py}  & \texttt{irrep\_identification.json} \\
                                          & \texttt{sim\_p\_polynomial\_structure.py} & \texttt{p\_polynomial\_results.json} \\
\ref{thm:closure-spec} (Closure spec)    & \texttt{sim\_closure\_sigma\_action.py} & \texttt{closure\_sigma\_action\_results.json} \\
\ref{thm:univ-rep} (Generation)          & \texttt{sim\_generation\_operator.py}   & \texttt{generation\_operator\_results.json} \\
\ref{thm:sigma-pair} ($\sigma$-pair)     & \texttt{sim\_closure\_sigma\_action.py} & \texttt{closure\_sigma\_action\_results.json} \\
\ref{thm:triad-iso} (Triad iso)          & \texttt{sim\_triad\_isomorphism.py}     & \texttt{triad\_isomorphism\_results.json} \\
\ref{thm:rpi} ($r(\pi)$, maximal order)   & \texttt{sim\_theta\_maximal\_order.py}  & \texttt{theta\_maximal\_order\_results.json} \\
\ref{thm:rpi} (Lipschitz proxy, odd $\pi$) & \texttt{sim\_representation\_numbers.py} & \texttt{representation\_numbers\_results.json} \\
\ref{thm:mult} (Multiplicativity)        & \texttt{sim\_multiplicativity.py}       & \texttt{multiplicativity\_results.json} \\
\ref{thm:euler-bridge} (L-function, exact) & \texttt{sim\_theta\_maximal\_order.py}  & \texttt{theta\_maximal\_order\_results.json} \\
\hline
\end{tabular}
\end{center}

Each simulation runs in well under a minute on a standard laptop and passes its
acceptance criteria. The maximal-order representation counts and the exact
identity $L(\Theta_\II,s)=\zeta_K(s)\zeta_K(s-1)$ are verified by complete
short-vector enumeration in \texttt{sim\_theta\_maximal\_order.py}; the earlier
\texttt{sim\_theta\_L\_function.py} (which summed the divisor-sum Dirichlet
series rather than the geometric theta, and the Lipschitz-suborder
\texttt{sim\_representation\_numbers.py}) are retained for provenance. Full
output is in \texttt{repro/output/}.

\bigskip

\bibliographystyle{plainnat}
\begin{thebibliography}{9}

\bibitem{rhFormalSmart}
L.~Smart,
\emph{The Riemann Hypothesis as a $\sigma$-fixed spectral statement on $V_{600}$ — a $(\delta)$-formal reformulation},
Open research preprint, Institute of Vibrational Field Dynamics, 2026.

\bibitem{existenceClosure}
L.~Smart,
\emph{Existence as Closure},
Open research preprint, Institute of Vibrational Field Dynamics, 2026.

\bibitem{lifeAsClosure}
L.~Smart,
\emph{Life as Closure},
Open research preprint, Institute of Vibrational Field Dynamics, 2026.

\bibitem{spectralBridge24600}
L.~Smart,
\emph{From Schl\"afli Decomposition to Spectral Bridge: A First Closure-Projection Channel in $V_{600}$},
Open research preprint, Institute of Vibrational Field Dynamics, 2026.

\bibitem{conwaySloane}
J.~H.~Conway and N.~J.~A.~Sloane,
\emph{Sphere Packings, Lattices and Groups},
Springer-Verlag, 3rd ed., 1999.

\bibitem{brouwerCohenNeumaier}
A.~E.~Brouwer, A.~M.~Cohen, and A.~Neumaier,
\emph{Distance-Regular Graphs},
Springer-Verlag, 1989.

\bibitem{eichler1958}
M.~Eichler,
\emph{Quadratische Formen und Modulfunktionen},
Acta Arithm., 4 (1958), 217--239.

\bibitem{hijikata1974}
H.~Hijikata,
\emph{Explicit formula of the traces of Hecke operators for $\Gamma_0(N)$},
J.~Math.~Soc.~Japan, 26 (1974), 56--82.

\bibitem{vigneras}
M.-F.~Vign\'eras,
\emph{Arithm\'etique des Alg\`ebres de Quaternions},
Lecture Notes in Mathematics 800, Springer-Verlag, 1980.

\end{thebibliography}

\end{document}
